% ============================================================================
%  CHAPITRE 1 -- FONDEMENTS ET ÉTAT DE L'ART
% ============================================================================
\chapter{Fondements et état de l'art}
\label{chap:contexte}

Ce chapitre n'introduit que les notions utilisées par la suite : le théorème de séparation et ses limites, les
modèles de canal retenus, la chaîne numérique classique et ses mécanismes d'erreur, puis les travaux sur
lesquels s'appuie la communication sémantique apprise.

\section{Le théorème de séparation et sa limite pratique}
\label{sec:separation}

Shannon montre qu'un système de communication peut être décomposé sans perte d'optimalité en deux blocs
indépendants~\cite{shannon1948} : un \textbf{codage de source}, qui retire la redondance de l'information, et un
\textbf{codage de canal}, qui réintroduit une redondance calculée pour résister au bruit. C'est le fondement de
toutes les architectures déployées : JPEG puis Wi-Fi, H.264 puis 4G.

Ce résultat est cependant valable à la limite des blocs de longueur infinie, pour un canal stationnaire dont
la statistique est connue de l'émetteur. Trois hypothèses tombent en pratique :

\begin{itemize}
  \item \textbf{Les blocs sont courts.} Une image CIFAR-10 au point de fonctionnement étudié dispose de
        \num{256} symboles complexes, soit quelques centaines de bits utiles. Le gain asymptotique des codes
        correcteurs n'est pas atteint à cette longueur.
  \item \textbf{Le SNR réel diffère du SNR prévu.} L'émetteur choisit son schéma de modulation et de codage
        avant de connaître l'état exact du canal.
  \item \textbf{Le critère optimisé n'est pas celui du destinataire.} La chaîne minimise une distorsion de
        reconstruction, alors que la tâche finale est souvent une reconnaissance.
\end{itemize}

\begin{definition}[Effet falaise]
Comportement en tout-ou-rien d'une chaîne numérique conçue pour un SNR cible : au-dessus du seuil, la qualité
sature et la marge disponible n'est pas exploitée ; en dessous, le décodeur de canal décroche, le train binaire
devient faux et l'information source est perdue en bloc plutôt que dégradée progressivement.
\end{definition}

\section{Modèles de canal}
\label{sec:canaux}

Le rapport signal sur bruit s'exprime en décibels, $\mathrm{SNR_{dB}} = 10\log_{10}(P_{\text{signal}} /
P_{\text{bruit}})$. Une valeur de \SI{0}{\decibel} signifie autant de bruit que de signal, une valeur de
\SI{20}{\decibel} cent fois moins de bruit que de signal. Deux modèles sont utilisés dans ce travail.

\paragraph{Canal AWGN.} Le modèle thermique idéal : à chaque symbole complexe émis s'ajoute un bruit gaussien
indépendant,
\begin{equation}
  y = x + n, \qquad n \sim \mathcal{CN}(0,\sigma^2), \qquad \sigma^2 = 10^{-\mathrm{SNR_{dB}}/10} .
  \label{eq:awgn}
\end{equation}

\paragraph{Canal de Rayleigh par blocs.} Modèle d'une liaison mobile : le signal atteint le récepteur par
plusieurs trajets réfléchis dont les phases s'additionnent, de sorte que l'amplitude reçue devient aléatoire,
\begin{equation}
  y = h\,x + n, \qquad h \sim \mathcal{CN}(0,1) .
  \label{eq:rayleigh}
\end{equation}
Le coefficient $h$ est tiré une fois par image (d'où « par blocs ») et reste constant pendant toute sa
transmission. Le récepteur est supposé connaître $h$ (CSI parfaite) et divise par $h$ : cette
\textbf{égalisation par forçage à zéro} ramène le problème à un canal AWGN de rapport signal sur bruit
instantané $\mathrm{SNR}\cdot|h|^2$, où $|h|^2$ suit une loi exponentielle de moyenne 1. Lorsque le tirage est
défavorable, l'évanouissement est profond et toute l'image en souffre.

\paragraph{Contrainte de puissance.} Un émetteur ne peut pas augmenter librement son amplitude pour couvrir le
bruit. La puissance moyenne est donc normalisée à 1 par symbole complexe. Sans cette contrainte, un système
appris obtiendrait un gain trivial en amplifiant sa sortie, et la comparaison avec la chaîne numérique
perdrait tout sens.

\section{La chaîne numérique classique}
\label{sec:chaine_classique}

\paragraph{Modulation QAM.} Les bits sont groupés et associés aux points d'une constellation du plan complexe :
QPSK transporte 2~bits par symbole, la 16-QAM 4~bits, la 64-QAM 6~bits. Plus la constellation est dense, plus
le débit est élevé, mais plus les points sont proches et donc sensibles au bruit. Le \textbf{codage de Gray}
numérote les points de sorte qu'une confusion entre voisins immédiats ne coûte qu'un seul bit faux.

\paragraph{Code convolutif et poinçonnage.} Le code retenu est le code convolutif de contrainte $K=7$ et de
polynômes générateurs $(133,171)_8$, celui du Wi-Fi 802.11a~\cite{ieee80211a}. Chaque bit émis dépend des sept
derniers bits d'entrée. Le rendement natif $1/2$ est porté à $2/3$ ou $3/4$ par \textbf{poinçonnage}, c'est-à-dire
en supprimant certains bits de redondance selon un motif fixe.

\paragraph{Décodage souple et Viterbi.} Le démodulateur ne tranche pas entre 0 et 1 : pour chaque bit il calcule
un log-rapport de vraisemblance $\mathrm{LLR} = \log\big(P(0)/P(1)\big)$, dont le signe indique le bit le plus
probable et l'amplitude la confiance. L'algorithme de Viterbi~\cite{viterbi1967} recherche ensuite le chemin le
plus vraisemblable dans le treillis à $2^{K-1} = 64$ états. Tant que le canal reste au-dessus du seuil du code,
la correction est quasi parfaite ; en dessous, le décodeur choisit le mauvais chemin et les erreurs passent en
masse. Le couple (modulation, rendement) est appelé \textbf{MCS}.

\section{Travaux antérieurs}
\label{sec:etat_art}

\paragraph{Communication sémantique.} Le domaine a été reformulé récemment autour de l'idée de Weaver : ne
transmettre que ce qui modifie l'interprétation du destinataire~\cite{qin2021,gunduz2023}. Xie
\emph{et al.}~\cite{xie2021} ont proposé un système de bout en bout pour le texte, Zhang
\emph{et al.}~\cite{zhang2023} une variante dont l'émetteur ignore la tâche du récepteur, et Fu
\emph{et al.}~\cite{fu2023} une quantification vectorielle du latent.

\paragraph{JSCC profond pour l'image.} Bourtsoulatze \emph{et al.}~\cite{bourtsoulatze2019} ont introduit le
codage conjoint source-canal convolutif dont notre système syntaxique est la reprise : encodeur CNN, symboles
analogiques, canal inséré comme une couche du réseau, décodeur CNN. Kurka et Gündüz~\cite{kurka2021} ont montré
qu'un seul modèle peut couvrir plusieurs rapports de bande, propriété que nous exploitons par troncature des
canaux latents. Xu \emph{et al.}~\cite{xu2022} y ont ajouté des modules d'attention, Dai
\emph{et al.}~\cite{dai2022} une transformée non linéaire.

\paragraph{Vidéo et transmission orientée tâche.} DeepWiVe~\cite{tung2022} et la transmission sémantique
vidéo~\cite{wang2023} étendent l'approche aux séquences. Shao \emph{et al.}~\cite{shao2022} formulent la
transmission orientée tâche comme un goulot d'étranglement informationnel : ne transmettre que ce dont
l'inférence a besoin. Notre extension du chapitre~\ref{chap:extensions} en est une version simple et directe.

\paragraph{Pertes perceptuelles.} Le terme de caractéristiques de notre perte sémantique est une perte
perceptuelle au sens de Johnson \emph{et al.}~\cite{johnson2016} : comparer les activations internes d'un
réseau gelé plutôt que les pixels. La nouveauté n'est pas la perte elle-même, mais son usage dans une chaîne de
transmission et surtout le protocole de mesure qui l'accompagne (chapitre~\ref{chap:protocole}).

\begin{table}[H]
\centering
\caption{Positionnement des trois systèmes comparés dans ce projet}
\label{tab:positionnement}
\small
\begin{tabularx}{\textwidth}{@{}p{3.2cm}cc>{\raggedright\arraybackslash}X@{}}
\toprule
\textbf{Système} & \makecell{\textbf{Bits}\\\textbf{explicites}} & \makecell{\textbf{Effet}\\\textbf{falaise}} & \textbf{Critère optimisé} \\
\midrule
Chaîne séparée & \cmark & \cmark & distorsion de source, puis fiabilité binaire \\
JSCC syntaxique & \xmark & \xmark & erreur quadratique sur les pixels \\
JSCC sémantique & \xmark & \xmark & pixels \emph{et} caractéristiques \emph{et} tâche \\
\bottomrule
\end{tabularx}
\end{table}
